% Revised 22 October 2026: completed final system and bounded future work.
\chapter{Conclusion and Future Work}
\label{ch:conclusion}

\section{Conclusion}
\label{sec:concl}

This thesis presents and evaluates a compact student for multi-person human mesh recovery with Anny. The model
combines a ViT-B encoder, a learned feature projector and prediction heads initialized from a ViT-L teacher.
Training combines AnnyOne and converted 3DPW supervision with feature and output distillation. At inference
time only the student and Anny decoding layer are required.

The conversion pipeline bridges SMPL-annotated real images and Anny-space supervision through geometric
joint correspondence, virtual torso joints and a shared shape per actor. Independent checks document its
accuracy and annotation limits. Consistent head anchors and fixed output corrections form part of the
teacher/student configuration, with supporting analyses provided in the appendices.

The student has \StudentParams\,M parameters and \StudentGFLOPs\,GFLOPs in the recorded configuration,
compared with \FinalParams\,M and \FinalGFLOPs\,GFLOPs for the teacher. Paired L40S forward latency is
\StudentMsLFortyS\,ms versus \TeacherMsLFortyS\,ms. Student PVE is \StudentPVE\,mm on AnnyOne and
\StudentThreeDPW\,mm on converted 3DPW. The teacher remains more accurate on both sets. The student improves
over the released baseline's \NAVERThreeDPW\,mm 3DPW PVE, although real-image training exposure differs.
These results establish a compactness and speed advantage with an accuracy trade-off.

The video application processes identical source frames with the student and baseline and saves their meshes
beside the input. In the completed \VideoFrames{}-frame walking example the student has \VideoStudentMatched{}
tracked predictions and the baseline \VideoBaselineMatched{}. This provides an initial video demonstration;
a temporal-quality claim requires inspection of the saved videos and measurements. The implementation remains
a frame-wise image model with tracking, rather than a learned temporal predictor.

\section{Future Work}
\label{sec:future}

\paragraph{Accuracy at reduced cost.}
Longer training, distillation-weight ablations and alternative feature transfer strategies could reduce the
student's error while retaining its lower cost. They need matched evaluation against the current student,
teacher and a baseline adapted to the same real-image data. This report uses the verified 30-epoch student;
a later run should replace it only after its checkpoint and complete evaluation are available.

\paragraph{Temporal evaluation and modelling.}
The next step is to inspect the walking overlays and parameter variation, then examine clips with occlusion
and camera motion. Synchronized inference FPS should be reported with GPU, precision and resolution.
If substantial jitter is observed, track-wise shape aggregation, smoothing or learned temporal refinement can
be evaluated against the unsmoothed baseline. Stability and reconstruction accuracy should be measured together.

\paragraph{Real-image shape.}
Age-group labels, body measurements or additional real-image Anny fits could improve supervision beyond
constant converter values. Evaluation should include children and distinguish the body model's ability to
generate child proportions from the network's ability to infer them reliably from real photographs.

\paragraph{Reliability and reproducibility.}
Repeated runs, uncertainty estimates and a larger held-out evaluation would strengthen comparisons.
Portable configuration and a documented release process would make the system easier to verify.
These are possible extensions, rather than completed experiments.
